\section{导读：机器人为什么既近又远}

本节先说明这期的阅读方式。谭杰是 Google DeepMind Robotics 团队高级研究科学家和技术负责人，他的视角很适合拆开机器人行业里常被混在一起的几个问题：仿真如何改变真实控制，强化学习解决了什么，大模型给机器人补了什么，机器人基座模型是否已经独立成学科，Gemini Robotics 1.5 的 Motion Transfer 在解决什么，以及为什么高质量数据仍是最硬的瓶颈。

本期最重要的平衡感是：机器人进展很快，但离大规模落地还远。谭杰一方面认为硅谷普遍把机器人视为即将发生的重大变革，也承认过去一年模型泛化有意外进展；另一方面他反复提醒，网上最好的 demo 往往不是代表性能力，现实落地从论文、live demo 到生产系统可能隔着多年距离。这个平衡让本期非常适合作为机器人领域的“冷静乐观”笔记。

\begin{importantbox}{本期核心命题}
机器人不是单一模型问题，而是大脑、小脑、本体、数据、仿真、触觉、安全和落地场景共同作用的综合问题。大模型带来语言理解和 common sense，强化学习解决部分底层控制，世界模型和合成数据试图补足泛化，但高质量真实/仿真数据仍是核心瓶颈。
\end{importantbox}

\begin{knowledgebox}{视觉策略说明}
本视频是固定访谈画面，没有 slides、白板或产品演示。正文只使用封面作为来源识别，正文图像全部为自制概念图，用来解释 Sim2Real、VLA、Gemini Robotics 1.5、Synthetic Data、世界模型 V-L-V、触觉和中美机器人分工。
\end{knowledgebox}

\subsection{本章小结}

EP121 的主线不是“机器人马上进入家庭”，而是解释为什么机器人领域已经出现范式信号、为什么落地仍慢、以及哪些问题决定接下来五到十年的路线。

